\section{Compilation: flattening, scheduling and backend-aware optimisation}
\label{sec:compiler}

The comparisons of this survey are possible because one source program is lowered
to every paradigm rather than being written once per paradigm. That lowering is a
compiler, and it is a contribution in its own right: the object it optimises is
not a machine but a \emph{protocol}, whose cost model is unlike a processor's, and
several standard compiler transformations change sign under it. This section
develops the intermediate representations, states what the optimiser is permitted
to do, and proves the two structural facts that distinguish the problem from
classical instruction selection.

\subsection{Two intermediate representations, and why there are two}
\label{sec:compiler:ir}

\begin{definition}
\label{def:flatop}
A \ib{flat program} over a set of wires $\mathcal{W}$ is a finite sequence of
instructions in static single-assignment form,
$\iota_1,\dots,\iota_m$, where each $\iota_j$ has an operation, a list of source
wires and a destination wire $d_j$, subject to: every $d_j$ is distinct, and every
source wire either is an input or equals $d_{j'}$ for some $j' < j$. We write
$\mathcal{P}$ for the set of flat programs. The operations are constants, ring
addition and subtraction, multiplication, signed and unsigned comparison,
equality, bitwise operations, shifts, selection and, in the non-deterministic
setting, division.
\end{definition}

Flat programs are the representation a proof system wants, because a constraint
system is a conjunction of relations over named wires and SSA gives exactly that.
They are also what a homomorphic evaluator wants, because a homomorphic evaluator
is an interpreter over a dataflow graph. They are, however, the wrong
representation for costing a secret-sharing protocol, for three reasons that
Definition~\ref{def:bytecode} then addresses.

\begin{definition}
\label{def:bytecode}
A \ib{vectorised bytecode program} is a finite sequence of instructions over
registers, where each register carries a \ib{width} $\omega \in \N$, a
\ib{domain} $\mathcal{D} \in \{\Z_{2^k},\ \mathbb{F}_p,\ \mathbb{F}_2\}$ and a
\ib{secrecy} tag in $\{\textsf{clear},\textsf{secret}\}$. Instructions are
partitioned into \ib{local} ones, which consume no communication, and
\ib{interactive} ones, which consume one round and possibly correlated
randomness. An instruction of width $\omega$ denotes $\omega$ independent
applications of its operation.
\end{definition}

The three deficiencies of Definition~\ref{def:flatop} for protocol costing are now
visible. It is scalar, so a dot product of $784$ terms appears as $784$
multiplications not as as one instruction of width $784$, and a cost model
reading it would charge $784$ rounds where a protocol pays one. It carries no
domain, so it cannot express that the same multiplication costs one Beaver triple
over $\Z_{2^k}$ and roughly $k$ AND gates over $\mathbb{F}_2$. And it has no
notion of an \emph{opening}, which is the unit of communication in a
secret-sharing protocol and does not correspond to any single arithmetic
operation.

\begin{definition}
\label{def:lowering}
A \ib{lowering} is a map $\mathcal{L} : \mathcal{P} \to \mathcal{B}$ from flat
programs to bytecode programs. It is \ib{faithful} when for every $P$ and every
input assignment, the outputs of $P$ and $\mathcal{L}(P)$ agree under the
plaintext semantics of both.
\end{definition}

Our compiler keeps both representations and one faithful lowering between them,
in place of a single universal intermediate representation. The reason is not
expedience: the two have genuinely different notions of cost, and a
representation that served both would have to carry the union of their
annotations while permitting only the transformations sound for both. Keeping the
flat form as the input to lowering, and confining protocol-aware reasoning to the
bytecode, means that a rewrite justified by a protocol's cost model cannot
accidentally be applied where the paradigm makes it unsound.

\subsection{Flattening and the classical passes}
\label{sec:compiler:flatten}

Flattening turns a nested expression into a flat program. Its interesting
properties are not the traversal, which is standard, but the interaction of the
classical passes with the cost model that follows.

\begin{definition}
\label{def:passes}
Let $P$ be a flat program with designated output wire $r$.
\ib{Constant folding} replaces any instruction all of whose sources are constants
by a constant. \ib{Common-subexpression elimination} identifies instructions with
equal operation and equal source wires, retaining one and substituting its
destination for the others. \ib{Dead-wire elimination} removes instructions whose
destinations are not reachable from $r$ in the dependence graph.
\end{definition}

\begin{proposition}
\label{prop:passes-sound}
Constant folding, common-subexpression elimination and dead-wire elimination
preserve the value computed at $r$, and dead-wire elimination applied after
common-subexpression elimination removes at least as many instructions as either
order of the two applied once.
\end{proposition}

\begin{proof}
Soundness of folding is that the operations are functions of their arguments;
soundness of elimination is that equal operations on equal arguments yield equal
results, so substituting one destination for another preserves every later value;
soundness of dead-wire removal is that SSA destinations are distinct, so an
instruction whose destination is unreachable from $r$ cannot influence $r$. For
the ordering claim, note that eliminating a common subexpression can only decrease
the set of reachable wires: it replaces occurrences of one destination by another
and never introduces a new dependence. Hence the set of dead wires after
elimination contains the set of dead wires before it, and the second pass removes
at least as many.
\end{proof}

Proposition~\ref{prop:passes-sound} justifies the pass order our compiler uses. It
is worth stating because the intuition from classical compilation---that
elimination and dead-code removal are independent---is false here in a useful
direction: collapsing duplicates is what \emph{creates} the dead wires, so running
the removal afterwards is strictly better and not merely tidier.

\subsection{Scheduling: rounds as a graph height}
\label{sec:compiler:sched}

\begin{definition}
\label{def:schedule}
A \ib{schedule} of a bytecode program assigns to each instruction a level in $\N$
such that every instruction's level exceeds the level of each instruction
producing one of its sources when that instruction is interactive, and is at least
that level when it is local. The \ib{as-soon-as-possible} schedule assigns each
instruction the least level satisfying these constraints.
\end{definition}

Proposition~\ref{prop:asap} established that the as-soon-as-possible schedule's
number of interactive levels equals the maximum number of interactive operations
on a directed path, and that no schedule does better. Two corollaries drive the
optimiser.

\begin{corolary}
\label{cor:batching}
Under the as-soon-as-possible schedule, all interactive instructions at one level
are mutually independent and may be issued as a single batched protocol
invocation. Consequently a layer of $n$ independent comparisons costs the rounds
of one comparison, and the round count of a network is a function of its depth
alone.
\end{corolary}

\begin{proof}
Two instructions at the same level cannot be joined by a directed path, since a
path through an interactive instruction strictly increases the level. Independence
permits batching by the definition of the protocols involved. The consequence for
a network follows because the nonlinearities of one layer depend only on that
layer's affine map and not on each other.
\end{proof}

\begin{corolary}
\label{cor:reassoc}
Let $\circ$ be an associative operation whose bytecode realisation is interactive,
and let $e$ be a left-nested chain $x_1 \circ (x_2 \circ (\cdots \circ x_n))$.
Rebalancing $e$ into a tree of depth $\ceil{\log n}$ leaves the instruction count
unchanged and reduces the interactive depth from $n-1$ to $\ceil{\log n}$. If the
realisation of $\circ$ is \emph{local}, rebalancing changes nothing.
\end{corolary}

\begin{proof}
A binary tree with $n$ leaves has $n-1$ internal nodes regardless of shape, so the
instruction count is invariant; the height of a balanced tree is $\ceil{\log n}$
and of a chain is $n-1$, and by Proposition~\ref{prop:asap} the interactive height
is the round count. If $\circ$ is local it contributes no interactive levels in
either shape.
\end{proof}

Corollary~\ref{cor:reassoc} is the cleanest instance of the phenomenon this
section exists to formalise: reassociation is a transformation that changes
\emph{only} depth, leaving every other count identical, so its value is determined
entirely by whether the target protocol charges for depth. Over an arithmetic
sharing, addition is local and multiplication is interactive, so rebalancing a
product chain pays and rebalancing a sum chain does not. Over a binary protocol an
adder is itself a circuit with AND-depth, so rebalancing a sum chain pays as well.
Over a constant-round protocol neither pays. One transformation, three answers,
determined by the cost model and not by a heuristic.

\subsection{Backend-aware rewriting}
\label{sec:compiler:rewrite}

\begin{definition}
\label{def:rewrite}
A \ib{rewrite rule} is a partial map $\rho$ on bytecode programs that preserves
plaintext semantics. A \ib{cost model} for a protocol $\pi$ and deployment
$\delta$ is a function $\mathrm{cost}_{\pi,\delta} : \mathcal{B} \to \Rr_{\ge0}$.
A rule $\rho$ is \ib{profitable} for $(\pi,\delta)$ at $P$ when
$\mathrm{cost}_{\pi,\delta}(\rho(P)) < \mathrm{cost}_{\pi,\delta}(P)$.
\end{definition}

\begin{definition}
\label{def:baopt}
The \ib{backend-aware optimiser} takes a program $P$, a protocol $\pi$ and a
deployment $\delta$, and returns
$\rho_{i_k} \circ \cdots \circ \rho_{i_1}(P)$ obtained by considering candidate
rules and applying each only when it is profitable for $(\pi,\delta)$ at the
current program, with semantic preservation checked by differential evaluation
against the original on random inputs.
\end{definition}

The definition is greedy and makes no optimality claim; we state its limitations
below not leave them implicit. What it does provide is that the target
protocol is a \emph{parameter} of optimisation rather than a fixed assumption, and
the consequence is that the optimiser emits different circuits for different
protocols from the same source. We give the phenomenon a statement, since it is
the contribution.

\begin{proposition}
\label{prop:rewrite-sign}
There exist a rewrite rule $\rho$, a program $P$ and protocols $\pi_1,\pi_2$ such
that $\rho$ is profitable for $\pi_1$ at $P$ and unprofitable for $\pi_2$ at $P$.
\end{proposition}

\begin{proof}
Take $\rho$ to be strength reduction, replacing multiplication by a public
constant $c$ with a sum of shifts, and let $P$ multiply a secret by $c$. Under a
binary protocol, a $k\times k$ multiplier is $\Theta(k^2)$ AND gates while an
adder is $\Theta(k)$, so replacing a general multiplication by a small number of
shifts and additions reduces the gate count and hence the cost; $\rho$ is
profitable. Under an arithmetic sharing over $\Z_{2^k}$, multiplication by a
public constant is a local scaling of each share and costs nothing, whereas the
rewritten form is a sequence of local operations that is no cheaper and adds
instructions; $\rho$ is not profitable. The two protocols therefore assign
opposite signs to the same rule at the same program.
\end{proof}

\begin{remark}
\label{rem:persite}
Proposition~\ref{prop:rewrite-sign} concerns one rule at one program. A subtler
point is that profitability varies \emph{within} a program: on
$(x \cdot 2^a) + (y \cdot c)$ with $c$ not a power of two, strength reduction is
profitable at the first site and not at the second under a protocol where the
shift is free and the general constant is not. An optimiser that decides a rule
once for the whole program must therefore take both or neither, and will be wrong
at one site. Deciding per site costs a full cost evaluation per candidate, which
is affordable because evaluating the model is cheap, and it is what makes the
emitted circuits differ between protocols by a \emph{subset} relation rather than
arbitrarily.
\end{remark}

\begin{remark}
\label{rem:opt-limits}
The optimiser of Definition~\ref{def:baopt} is greedy with a fixed rule order and
evaluates candidate sites independently. It therefore misses any pair of rewrites
that is profitable only jointly, and it offers no approximation guarantee. We
state this not characterise the loss, because we have not measured it; the
honest claim is that the optimiser improves the cost model's prediction relative
to no optimisation and relative to a backend-blind choice, not that it is optimal.
\end{remark}

\subsection{Protocol assignment, and why the objective is not separable}
\label{sec:compiler:assign}

The final transformation chooses \emph{which} protocol evaluates each part of a
program, paying for conversions where the choice changes.

\begin{definition}
\label{def:assignment}
Let $\Sigma$ be a finite set of protocols and let $K$ be the set of operation
kinds occurring in $P$. An \ib{assignment} is a map $a : K \to \Sigma$. A
\ib{crossing} is an edge of the dependence graph whose endpoints are assigned
different protocols; realising it requires a \ib{conversion} instruction. The
\ib{assignment problem} is to minimise
$\mathrm{cost}_{\delta}\bigl(a\bigr)$, the cost of the program under $a$ with
conversions inserted at every crossing.
\end{definition}

\begin{proposition}
\label{prop:assign-nonsep}
The objective of Definition~\ref{def:assignment} is not separable over operations:
there are programs and deployments for which the cost of an assignment is not
equal to any sum of per-operation terms, even allowing pairwise terms for
crossings. Two mechanisms suffice, and either alone is enough:
\begin{enumerate}
  \item a level's latency is the \emph{maximum} over the interactive instructions
  it contains, not their sum;
  \item conversions at the same level are \emph{batched}, so the marginal cost of a
  crossing depends on whether another crossing already occurs at that level.
\end{enumerate}
\end{proposition}

\begin{proof}
For (1), consider a deployment with latency $L$ and unbounded bandwidth, and a
program consisting of $m$ mutually independent interactive operations. By
Proposition~\ref{prop:asap} the schedule has one interactive level, so the latency
contribution is $L$ regardless of $m$; any separable objective assigning a
positive term $c > 0$ to each such operation would report $mc$, which is not
constant in $m$. Moving one of the $m$ operations to another protocol changes the
cost by zero if that protocol's level is already occupied and by $L$ if it is not,
so the marginal cost of an individual assignment decision depends on the rest of
the assignment. For (2), let two crossings occur at one level; their combined
conversion cost is that of a single batched conversion, which is strictly less than
twice the cost of one whenever the conversion has a fixed component. Then the cost
of the assignment that makes both crossings is less than the sum of the costs of
the two assignments each making one, contradicting separability with pairwise
terms.
\end{proof}

\begin{corolary}
\label{cor:no-ilp}
An integer-linear formulation of the assignment problem with per-operation and
per-crossing variables and a linear objective cannot express the cost of
Definition~\ref{def:assignment} exactly.
\end{corolary}

\begin{proof}
A linear objective over such variables is by construction a sum of per-operation
and per-crossing terms, which Proposition~\ref{prop:assign-nonsep} shows is not
equal to the cost function.
\end{proof}

Corollary~\ref{cor:no-ilp} places our approach with respect to the prior work on
protocol assignment, which formulates the problem as an integer program and
therefore commits to a linear, separable cost. That formulation is a good
approximation in a regime where each operation's cost is dominated by bandwidth,
which is separable, and it becomes inaccurate in the regime where latency
dominates---which is the wide-area regime and precisely the one in which the
paradigms differ most. Because evaluating our cost model on a whole program is
inexpensive, we can afford to optimise the exact objective by search, giving up
optimality guarantees instead of fidelity. A learned surrogate for the cost model
would be the alternative if the search space grew beyond exact evaluation, and it
is the direction the recent literature takes; we note it as complementary and not as as competing.

\begin{remark}
\label{rem:assign-scope}
Our assignment is per operation \emph{kind}, not per instruction. The space is
therefore $\abs{\Sigma}^{\abs{K}}$ with $\abs{K}$ a few dozen, in place of
$\abs{\Sigma}^{m}$ with $m$ the instruction count, and the resulting assignment is
interpretable: a statement such as ``comparisons moved to the binary protocol'' is
a statement about the cost model. Per-instruction assignment is strictly more
expressive and is the natural next step; we do not claim it.
\end{remark}

\subsection{Summary of the compilation pipeline}
\label{sec:compiler:summary}

Collecting the above, a source program traverses: parsing to an abstract syntax
tree; flattening to a flat program (Definition~\ref{def:flatop}) with constant
folding, common-subexpression elimination and dead-wire elimination
(Proposition~\ref{prop:passes-sound}); lowering to vectorised bytecode
(Definition~\ref{def:bytecode}) with domain and secrecy annotations; backend-aware
rewriting against the cost model of the chosen protocol
(Definition~\ref{def:baopt}); protocol assignment with conversions priced at
crossings (Definition~\ref{def:assignment}); as-soon-as-possible scheduling
(Definition~\ref{def:schedule}); and execution on a protocol engine or, for the
proof and homomorphic paradigms, direct interpretation of the flat program.

Two properties of this pipeline are what the rest of the survey relies on. Every
paradigm receives the \emph{same} source, so a comparison between paradigms is not
confounded by a difference in the program. And the cost model that guides
optimisation is the same one whose predictions we validate against measurement,
so a claim that an optimisation helps is falsifiable in the same terms as a claim
about a protocol.
